\chapter{Funciones y Ámbitos}
\label{chap:functions}
{\huge$\bigstar$}{\huge$\bigstar$} Intermedio

\section{Definición de Funciones}
\begin{lispcode}
(defun cuadrado (x) (* x x))
\end{lispcode}

\section{Parámetros}
\begin{lispcode}
(defun saludar (nombre &optional (saludo "Hola"))
  (format t "~A, ~A~%" saludo nombre))
(defun suma (&rest numeros) (reduce #'+ numeros))
(defun crear-punto (&key (x 0) (y 0)) (cons x y))
\end{lispcode}

\section{Funciones Anónimas}
\begin{lispcode}
(lambda (x) (* x x))
#'(lambda (x) (* x x))
\end{lispcode}

\section{Funciones Locales}
\begin{lispcode}
(flet ((doble (x) (* x 2)) (triple (x) (* x 3))) (+ (doble 5) (triple 5)))
(labels ((par? (n) (if (= n 0) t (impar? (1- n))))
         (impar? (n) (if (= n 0) nil (par? (1- n)))))
  (par? 4))
\end{lispcode}

\section{Closures}
\begin{lispcode}
(defun make-contador (&optional (inicial 0))
  (let ((cuenta inicial)) (lambda () (incf cuenta))))
(defvar *c* (make-contador 10))
(funcall *c*)
\end{lispcode}

\section{Múltiples Valores}
\begin{lispcode}
(values 10 20 30)
(multiple-value-bind (a b c) (values 1 2 3) (list a b c))
\end{lispcode}

\begin{examplebox}[title=Contador con estado encapsulado]
\begin{lispcode}
(defun make-acumulador ()
  (let ((total 0)) (lambda (n) (incf total n))))
\end{lispcode}
\end{examplebox}

\begin{examplebox}[title=Divmod con múltiples valores]
\begin{lispcode}
(defun divmod (a b) (values (floor a b) (mod a b)))
(multiple-value-bind (c r) (divmod 17 5) (list c r))
\end{lispcode}
\end{examplebox}

